% 原书第 84—88 页；第 6.5 节、第 6.6 节导言及第 6.6.1 节。
\section{关于进一步降维的考虑}
\label{sec:ch06-dimension-reduction}

当 $d$ 相当大时，$m=(2^d-1)n$ 个基函数无法存入内存，使得最小损失估计量 $\psi_{n,\lambda}$ 的计算不可行。由于经验风险 $P_nL(\psi)$ 只通过 $n$ 个值 $\{\psi(W_i):i=1,\ldots,n\}$ 依赖于 $\psi$，我们或许能够进一步减少基函数数目，同时仍达到经验风险的最小值。我们的定理证明，任何达到这一最小值的 $\psi_{n,\lambda}$，都会以所需的速率收敛到 $\psi_{0,\lambda}$。事实上，只需达到一个近似最小值，且近似误差小于这一速率即可。这表明，对于有限样本，也许只使用这些基函数的一个小得多的子集就足够了；随着样本量增加，再逐渐纳入所有类型的基函数，从而能够任意精确地逼近任何函数。开发在计算上可行、并能逼近所需 $\psi_{n,\lambda}$ 的算法，将是未来研究的一个重要方向。

我们提出以下策略，来定义一个纳入 HAL 估计量的超级学习器。关键步骤是构造一个嵌套的候选估计量序列，使序列中的最后一个估计量为完整的 HAL 估计量。例如，第一个估计量可以是只纳入 HAL 表示式中一元指示函数的 Lasso 估计量，而第 $k$ 个估计量则纳入直至 $k$ 阶的所有多元指示函数，$k=1,\ldots,d$。不过，更合理的做法是使用一个细得多的候选估计量序列，使内存存储需求和计算耗时沿着这一序列逐渐增加。例如，可以为所有多元指示基函数提出一种可能由数据自适应确定的排序，先排一元，再排二元，依此类推。这还需要分别为一元指示基函数、二元指示基函数等排序。然后，可以将候选估计量序列定义为：第 $j$ 个估计量是纳入这一基函数序列中前 $K_j$ 个基函数的 Lasso，其中 $K_1<K_2<\cdots<K_M$，$j=1,\ldots,M$，$K_M$ 是 HAL 估计量中的基函数总数。将这些候选估计量全部纳入超级学习器的算法库，根据预言不等式，这个超级学习器的表现至少与包含全部 $K_M$ 个基函数的完整 HAL 估计量一样好。实际计算时，可以不去求出完整的超级学习器，而是沿着这一序列依次计算候选估计量，每次记录交叉验证风险。一旦交叉验证风险趋于平稳，甚至开始恶化，就将最后计算的估计量作为最终估计量。这一方案的有效性依赖于如下假设：剩余序列中更激进的估计量，不会优于已选出的估计量。这样，对于给定样本量 $n$，我们预期所选基函数的数目能以 $O(n)$ 为界，从而使该估计量可以计算。

\section{应用}
\label{sec:ch06-applications}

前面将 HAL 介绍为一种一般的非参数估计量，用来估计定义为损失函数期望最小化解的 $d$ 元函数值参数。本节考虑将 HAL 用于实值结局的条件均值估计。具体而言，我们讨论观测数据由随机变量 $O=(W,Y)\sim P_0\in\mathcal M$ 的 $n$ 个独立同分布副本组成的情形，其中 $\mathcal M$ 是非参数统计模型。我们对这个模型施加的唯一约束是：对每个 $P\in\mathcal M$，由 $P$ 决定的、给定 $W$ 时 $Y$ 的条件均值，记作 $\bar Q_P$，都具有有限的变差范数。我们考虑用高度自适应 Lasso 估计 $\bar Q_0=\argmin_{\bar Q}P_0L(\bar Q)$，其中，$L(\bar Q)(o)=\{y-\bar Q(w)\}^2$ 为平方误差损失。

\subsection{高度自适应 Lasso 的构造}
\label{subsec:ch06-constructing-hal}

回顾一下，HAL 可以看作在一类特殊的指示基函数线性组合上，最小化经验风险所得的解，其约束是系数绝对值之和不大于某个由数据自适应选出的常数。本小节用简单的一元和二元情形，说明如何构造这些基函数和估计量。

考虑观测数据包含 $n=500$ 个独立副本，其中，$W\sim\operatorname{Uniform}(-4,4)$，$Y=2\sin(\pi/2|W|)+\epsilon$，而 $\epsilon$ 独立于 $W$，来自 $\operatorname{Normal}(0,1)$ 分布。HAL 使用的基函数，由观测数据取值处的 $n$ 个指示函数组成：$\phi_j(w)=I(w\geq w_j)$，$j=1,\ldots,n$。为选择变差范数的上界，我们使用十折交叉验证，从 100 个可能的上界中选择，其范围从 0 到约 350。图 6.1 展示了其中三个上界对应的拟合。实线是 HAL 估计量，使用交叉验证选出的值 $\lambda_n=13.9$。虚线和点线分别对应小于和大于真实变差范数 $\|\bar Q_0\|_v=16$ 的上界。图底部的短刻线，位于 $\bar Q_n$ 中系数非零的 46 个支撑点上。将变差范数上界取为 4.8（虚线），明显使数据被过度平滑；取为 35.2 则似乎能提供合理的逼近，与 HAL 估计量的预测相近。不过，这个更大的上界在支撑集边缘附近，似乎确实产生了更多噪声。根据理论，任何大于真实范数的上界，都会使所得估计量具有上一章所确立的性质。然而，在有限样本中，HAL 估计量允许选取小于真实范数的上界，因此会表现得更好。预言不等式保证：只要候选上界中至少有一个大于真实范数，我们最终就会选出一个大于真实变差范数的上界。
% 图内 lambda_n=13.4、正文 lambda_n=13.9，均照原书保留，不统一。

\begin{figure}[htbp]
 \centering
 \includegraphics[width=.88\textwidth]{figures/ch06_fig1_source.pdf}

 \caption{一元情形下的高度自适应 Lasso}
 \label{fig:ch06-hal-univariate}
\end{figure}
\FloatBarrier

下面考察二元情形，且 $W$ 含有一个离散分量。令 $W_1$ 来自 $\operatorname{Uniform}(-4,4)$ 分布，另独立地从 $\operatorname{Bernoulli}(0.5)$ 分布中抽取 $W_2$。设 $Y=-0.5W_1+W_2W_1^2/2.75+W_2+\epsilon$，其中，$\epsilon$ 来自 $\operatorname{Normal}(0,1)$ 分布。注意，这个数据生成分布意味着，在 $\bar Q_0$ 中，$W_1$ 与 $W_2$ 存在交互作用：当 $W_2=0$ 时，$W_1$ 与 $Y$ 的均值是线性关系，否则为二次关系。为在这一情形下构造 HAL 估计量，我们先创建 $n$ 个基函数，对应于 $W_1$ 观测值处的指示函数：$\phi_{1,j}(w)=I(w_1\geq w_{1,j})$，$j=1,\ldots,n$。接着，加入仅由 $W_2$ 构成的子集所对应的基函数：$\phi_{2,j}(w)=I(w_2\geq w_{2,j})$，$j=1,\ldots,n$。注意，由于 $W_2$ 是二值变量，只需加入一个不同的基函数，即 $\phi_2(w)=I(w_2\geq1)$。最后，创建如下形式的二元基函数：$\phi_{12,i}(w)=I(w_1\geq w_{1,j},w_2\geq w_{2,j})$，$j=1,\ldots,n$。由于 $W_2$ 是二值变量，这类基函数的数目小于 $n$。对于任何满足 $w_{2,j}=0$ 的 $j$，都不必再加入基函数 $\phi_{12,j}(w)$，因为通过纳入 $\bar\phi_{1,j}$，我们已经在相应的零边上放置了支撑点。这说明，给定样本中的基函数数目至多为 $n(2^d-1)$；实际数目则取决于具体数据集，可能少得多。

\begin{figure}[htbp]
 \centering
 \includegraphics[width=.88\textwidth]{figures/ch06_fig2_source.pdf}

 \caption{二元情形下的高度自适应 Lasso}
 \label{fig:ch06-hal-bivariate}
\end{figure}

图 6.2 展示了从这一数据生成机制中随机抽取的、样本量为 $n=500$ 的数据。两条曲线分别对应 $W_2=1$ 时 $\bar Q_0$ 的估计（上方虚线），以及 $W_2=0$ 时的估计（下方实线）。图底部的实线短刻线，表示 $\bar Q_n$ 中系数非零的一元基函数。因此，这些短刻线同时对应于 $\bar Q_n(\cdot,0)$ 和 $\bar Q_n(\cdot,1)$ 的跳跃。虚线短刻线表示系数非零的二元基函数，因此对应于 $\bar Q_n(\cdot,1)$ 跳跃而 $\bar Q_n(\cdot,0)$ 不跳跃的取值。注意，正如预期，这些短刻线最常出现在 $W_1>2$ 的位置：在这些取值处，$\bar Q_0(w_1,0)$ 随 $w_1$ 递减，而 $\bar Q_1(w_1,1)$ 递增。这个例子说明，HAL 估计量使用的高阶基函数，其作用类似于标准回归方法中的交叉乘积交互项。
% 原书 phi_{12,i} 与索引 j、bar phi_{1,j}、末句 bar Q_1 均按印刷保留。
\FloatBarrier
